-[9] Lớp 9

----[K] Khoa học tự nhiên

-------[0] Mở đầu

----------[1] Nhận biết một số dụng cụ, hoá chất. Thuyết trình một vấn đề khoa học

-------[1] Năng lượng cơ học

----------[1] Động năng. Thế năng

----------[2] Cơ năng

----------[3] Công và công suất

-------[2] Ánh sáng

----------[1] Khúc xạ ánh sáng
-------------[1] Hiện tượng khúc xạ và định luật khúc xạ ánh sáng
-------------[2] Chiết suất, mối quan hệ giữa góc tới và góc khúc xạ
-------------[3] Giải thích hiện tượng thực tế: đũa gãy trong nước, ảo ảnh sa mạc, đáy hồ nông hơn thực tế

----------[2] Phản xạ toàn phần

----------[3] Lăng kính

----------[4] Thấu kính
-------------[1] Thấu kính hội tụ và phân kì: đặc điểm, đường truyền tia sáng
-------------[2] Dựng ảnh qua thấu kính hội tụ
-------------[3] Công thức thấu kính (1/f = 1/d + 1/d')
-------------[4] Bài tập tính toán ảnh qua thấu kính
-------------[5] Ứng dụng thấu kính: kính lúp, kính hiển vi, kính thiên văn, máy ảnh

----------[5] Thực hành đo tiêu cự của thấu kính hội tụ

----------[6] Kính lúp. Bài tập thấu kính
-------------[1] Cấu tạo mắt, sự điều tiết, điểm cực cận, điểm cực viễn
-------------[2] Cận thị và cách khắc phục (kính phân kì)
-------------[3] Viễn thị và cách khắc phục (kính hội tụ)
-------------[4] Bảo vệ mắt: chế độ học tập, sử dụng thiết bị điện tử hợp lý

-------[3] Điện

----------[1] Điện trở. Định luật Ohm
-------------[1] Khái niệm điện trở, đơn vị đo
-------------[2] Định luật Ohm: I = U/R
-------------[3] Bài tập tính toán cơ bản theo định luật Ohm
-------------[4] Đặc tuyến Vôn-Ampe (đồ thị U-I)
-------------[5] Giải thích tại sao dây dẫn nóng lên khi có dòng điện chạy qua (bàn là, bếp điện...)
-------------[6] Điện trở phụ thuộc vào chiều dài, tiết diện, vật liệu dây dẫn
-------------[7] Điện trở suất của vật liệu
-------------[8] Bài tập tính điện trở theo công thức R = ρl/S
-------------[9] Giải thích tại sao dây tải điện dùng đồng/nhôm, dây mayso bếp dùng nikelin
-------------[A] Cấu tạo và hoạt động của biến trở
-------------[B] Bài tập biến trở trong mạch điện
-------------[C] Ứng dụng biến trở thực tế: núm chỉnh âm lượng, đèn mờ/sáng

----------[2] Đoạn mạch nối tiếp, song song
-------------[1] Đặc điểm đoạn mạch nối tiếp (I, U, R tương đương)
-------------[2] Đặc điểm đoạn mạch song song (I, U, R tương đương)
-------------[3] Bài tập tính toán mạch nối tiếp
-------------[4] Bài tập tính toán mạch song song
-------------[5] Bài tập mạch hỗn hợp (nối tiếp kết hợp song song)
-------------[6] Phân tích mạch điện trong gia đình: tại sao các thiết bị mắc song song

----------[3] Năng lượng của dòng điện và công suất điện
-------------[1] Khái niệm cường độ dòng điện, hiệu điện thế
-------------[2] Dụng cụ đo: Ampe kế, Vôn kế (cách mắc, đọc số liệu)
-------------[3] Thực hành đo cường độ dòng điện, hiệu điện thế trong mạch điện đơn giản
-------------[4] Ứng dụng: đọc hiểu thông số trên các thiết bị điện gia đình (V, A, W)
-------------[5] Công thức tính công suất điện P = UI
-------------[6] Công của dòng điện A = Pt = UIt
-------------[7] Bài tập tính công suất, công của dòng điện
-------------[8] Đọc hiểu nhãn thiết bị điện (220V - 1000W), tính cường độ dòng điện qua thiết bị
-------------[9] Tính toán tiền điện hàng tháng cho gia đình dựa trên công suất và thời gian sử dụng
-------------[A] Phát biểu và công thức định luật Joule-Lenz: Q = I²Rt
-------------[B] Bài tập tính nhiệt lượng tỏa ra trên điện trở
-------------[C] Giải thích nguyên lý hoạt động bàn là, bếp điện, ấm siêu tốc, bóng đèn dây tóc
-------------[D] Bài tập thực tế: tính thời gian đun nước bằng ấm siêu tốc (có hiệu suất)
-------------[E] Tác dụng của dòng điện đối với cơ thể người
-------------[F] Các quy tắc an toàn khi sử dụng điện
-------------[G] Cầu chì, aptomat (CB) - nguyên lý bảo vệ quá tải
-------------[H] Xử lý tình huống: người bị điện giật, chập điện, cháy do điện
-------------[I] Tiết kiệm điện năng trong gia đình (chọn thiết bị tiết kiệm, thói quen sử dụng)

-------[4] Điện từ

----------[1] Cảm ứng điện từ. Nguyên tắc tạo ra dòng điện xoay chiều
-------------[1] Từ trường của dòng điện trong dây dẫn thẳng, ống dây
-------------[2] Nam châm điện: cấu tạo, cách tăng lực từ
-------------[3] Ứng dụng nam châm điện: cần cẩu, rơ-le, chuông điện, khóa cửa từ
-------------[4] So sánh nam châm vĩnh cửu và nam châm điện
-------------[5] Lực từ tác dụng lên dây dẫn mang dòng điện đặt trong từ trường
-------------[6] Nguyên tắc hoạt động của động cơ điện
-------------[7] Ứng dụng động cơ điện: quạt, máy bơm, máy giặt, máy xay sinh tố
-------------[8] Giải thích hoạt động thiết bị điện trong nhà dùng động cơ điện
-------------[9] Hiện tượng cảm ứng điện từ (thí nghiệm Faraday)
-------------[A] Điều kiện xuất hiện dòng điện cảm ứng
-------------[B] Nguyên tắc hoạt động máy phát điện xoay chiều
-------------[C] Liên hệ thực tế: nhà máy thủy điện, phong điện, máy phát điện dự phòng

----------[2] Tác dụng của dòng điện xoay chiều
-------------[1] Đặc điểm dòng điện xoay chiều (tần số, chu kì)
-------------[2] Máy biến áp: cấu tạo, công thức U₁/U₂ = N₁/N₂
-------------[3] Vai trò máy biến áp trong truyền tải điện năng (giảm hao phí)
-------------[4] Giải thích hệ thống truyền tải điện: nhà máy → trạm biến áp → nhà dân

-------[5] Năng lượng với cuộc sống

----------[1] Vòng năng lượng trên Trái Đất. Năng lượng hoá thạch
-------------[1] Vai trò của năng lượng đối với sinh vật và con người
-------------[2] Các dạng năng lượng trong sinh hoạt và sản xuất
-------------[3] Phân tích chuỗi chuyển hóa năng lượng: Mặt Trời → quang hợp → thức ăn → năng lượng cơ thể

----------[2] Một số dạng năng lượng tái tạo
-------------[1] Năng lượng tái tạo và không tái tạo
-------------[2] Hiệu quả sử dụng năng lượng, giảm phát thải
-------------[3] Đánh giá thực trạng sử dụng năng lượng tại địa phương
-------------[4] Đề xuất giải pháp sử dụng năng lượng bền vững cho gia đình và cộng đồng

-------[6] Kim loại. Sự khác nhau cơ bản giữa phi kim và kim loại

----------[1] Tính chất chung của kim loại
-------------[1] Lý thuyết về tính chất vật lí của kim loại (dẻo, dẫn điện, dẫn nhiệt, ánh kim)
-------------[2] Lý thuyết về tính chất hóa học chung của kim loại (tác dụng với phi kim, acid, dung dịch muối)
-------------[3] Viết phương trình hóa học minh họa tính chất của kim loại
-------------[4] Bài tập nhận biết kim loại và hợp chất dựa vào tính chất đặc trưng
-------------[5] Bài tập định lượng về kim loại tác dụng với acid (HCl, H₂SO₄ loãng)
-------------[6] Bài tập định lượng về kim loại tác dụng với phi kim (O₂, Cl₂)
-------------[7] Bài tập liên quan đến ứng dụng thực tiễn của kim loại trong đời sống

----------[2] Dãy hoạt động hoá học
-------------[1] Lý thuyết về cấu tạo và ý nghĩa của dãy hoạt động hóa học
-------------[2] So sánh mức độ hoạt động của các kim loại
-------------[3] Xét điều kiện xảy ra phản ứng của kim loại với nước, acid, dung dịch muối
-------------[4] Viết phương trình hóa học của kim loại tác dụng với dung dịch muối
-------------[5] Bài tập nhận biết, tách chất, tinh chế dung dịch muối
-------------[6] Bài tập định lượng kim loại tác dụng với dung dịch muối (tăng giảm khối lượng)
-------------[7] Bài toán về kim loại kiềm, kiềm thổ (Na, K, Ca, Ba) tác dụng với dung dịch muối
-------------[8] Bài toán hỗn hợp kim loại tác dụng với dung dịch acid/muối

----------[3] Tách kim loại và việc sử dụng hợp kim
-------------[1] Lý thuyết về các phương pháp điều chế kim loại (thủy luyện, nhiệt luyện, điện phân)
-------------[2] Lý thuyết về hợp kim, gang, thép và ứng dụng
-------------[3] Viết phương trình hóa học điều chế kim loại từ hợp chất
-------------[4] Bài tập tách kim loại ra khỏi hỗn hợp
-------------[5] Bài toán định lượng liên quan đến quá trình nhiệt luyện (dùng CO, H₂ khử oxide kim loại)
-------------[6] Bài toán về thành phần phần trăm của hợp kim

----------[4] Sự khác nhau cơ bản giữa phi kim và kim loại
-------------[1] So sánh tính chất vật lí và hóa học cơ bản của kim loại và phi kim
-------------[2] Phân loại oxide (oxide acid, oxide base, oxide lưỡng tính, oxide trung tính)
-------------[3] Viết phương trình hóa học minh họa sự khác biệt giữa kim loại và phi kim
-------------[4] Bài tập nhận biết kim loại, phi kim và các hợp chất của chúng

-------[7] Giới thiệu về chất hữu cơ. Hydrocarbon và nguồn nhiên liệu

----------[1] Giới thiệu về hợp chất hữu cơ
-------------[1] Lý thuyết về hợp chất hữu cơ, hóa học hữu cơ, đặc điểm cấu tạo
-------------[2] Phân loại hợp chất hữu cơ (hydrocarbon, dẫn xuất hydrocarbon)
-------------[3] Viết công thức cấu tạo, công thức cấu tạo thu gọn của các hợp chất hữu cơ đơn giản
-------------[4] Tính phần trăm khối lượng các nguyên tố trong họp chất hữu cơ
-------------[5] Lý thuyết về đồng đẳng, đồng phân
-------------[6] Lập công thức phân tử khi biết phần trăm khối lượng các nguyên tố và khối lượng phân tử.
-------------[7] Lập công thức đơn giản của hợp chất hữu cơ
-------------[8] Lập công thức phân tử thông qua công thức đơn giản nhất.

----------[2] Alkane
-------------[1] Lý thuyết về dãy đồng đẳng, công thức chung, danh pháp, tính chất vật lí của alkane
-------------[2] Lý thuyết về tính chất hóa học đặc trưng (phản ứng thế, phản ứng cháy)
-------------[3] Viết phương trình phản ứng cháy của alkane
-------------[4] Viết các công thức cấu tạo và gọi tên một số alkane từ C1 đến C4
-------------[5] Bài toán tính phần trăm khối lượng mối nguyên tố trong hợp chất hữu cơ
-------------[6] Bài toán xác định công thức phân tử hợp chất hữu cơ khi biết phần trắm khối lượng và khối lượng phân tử
-------------[7] Lập công thức đơn giản của hợp chất hữu cơ
-------------[8] Lập công thức phân tử khi biết công thức đơn giản và các thông tin khác như khối lượng phân tử, số lượng nguyên tố khác,...

----------[3] Alkene
-------------[1] Lý thuyết về dãy đồng đẳng, công thức chung, danh pháp, tính chất vật lí của alkene
-------------[2] Lý thuyết về tính chất hóa học (phản ứng cộng, phản ứng trùng hợp, phản ứng oxi hóa)
-------------[3] Viết phương trình phản ứng cộng (H₂, Br₂, HX), trùng hợp, oxi hóa của alkene
-------------[4] Bài tập nhận biết alkene và alkane
-------------[5] Bài toán đốt cháy alkene
-------------[6] Bài toán alkene tác dụng với dung dịch Bromine
-------------[7] Bài toán tìm công thức phân tử alkene
-------------[8] Bài toán tổng hợp alkane và alkene

----------[4] Nguồn nhiên liệu
-------------[1] Lý thuyết về dầu mỏ, khí thiên nhiên, than đá và ứng dụng
-------------[2] Lý thuyết về chưng cất dầu mỏ, cracking và nhiên liệu sinh học
-------------[3] Bài toán tính nhiệt lượng tỏa ra khi đốt cháy nhiên liệu
-------------[4] Bài tập về các vấn đề môi trường liên quan đến việc khai thác và sử dụng nhiên liệu hóa thạch
-------------[5] Bài tập thực tiễn về nguồn nhiên liệu

-------[8] Ethylic alcohol và acetic acid

----------[1] Ethylic alcohol
-------------[1] Lý thuyết về cấu tạo, tính chất vật lí của ethylic alcohol
-------------[2] Lý thuyết về tính chất hóa học (phản ứng với Na, phản ứng cháy, phản ứng lên men)
-------------[3] Lý thuyết về độ cồn và cách pha chế
-------------[4] Bài toán định lượng liên quan đến phản ứng của ethylic alcohol với Na, phản ứng cháy, phản ứng lên men.
-------------[5] Bài toán về độ cồn và ứng dụng thực tiễn
-------------[6] Bài toán điều chế ethylic alcohol từ ethylene hoặc tinh bột và vận dụng thực tiễn

----------[2] Acetic acid
-------------[1] Lý thuyết về cấu tạo, tính chất vật lí của acetic acid
-------------[2] Lý thuyết về tính chất hóa học (tính acid, phản ứng este hóa)
-------------[3] Viết các phương trình hóa học minh họa tính chất
-------------[4] So sánh tính acid với các acid vô cơ và hữu cơ khác
-------------[5] Bài toán định lượng về tính acid (tác dụng với kim loại, base, oxide base, muối)
-------------[6] Bài toán về phản ứng este hóa (tính hiệu suất)
-------------[7] Bài toán về nồng độ dung dịch giấm ăn trong thực tế

-------[9] Lipid. Carbohydrate. Protein. Polymer

----------[1] Lipid
-------------[1] Lý thuyết về khái niệm, trạng thái tự nhiên và tính chất vật lí của lipid (chất béo)
-------------[2] Lý thuyết về tính chất hóa học (phản ứng thủy phân, phản ứng xà phòng hóa)
-------------[3] Viết phương trình thủy phân và xà phòng hóa chất béo
-------------[4] Bài tập về vai trò của chất béo và các vấn đề sức khỏe liên quan
-------------[5] Bài toán tính khối lượng xà phòng, glycerol từ chất béo (chỉ số acid, chỉ số xà phòng hóa)

----------[2] Carbohydrate. Glucose và saccharose
-------------[1] Lý thuyết về trạng thái tự nhiên, tính chất vật lí của glucose và saccharose
-------------[2] Lý thuyết về tính chất hóa học của glucose (phản ứng tráng gương, lên men) và saccharose (phản ứng thủy phân)
-------------[3] Viết các phương trình hóa học minh họa
-------------[4] Bài tập nhận biết glucose, saccharose và các carbohydrate khác
-------------[5] Bài toán về phản ứng tráng gương của glucose
-------------[6] Bài toán về phản ứng lên men rượu từ glucose
-------------[7] Bài toán về phản ứng thủy phân saccharose

----------[3] Tinh bột và cellulose
-------------[1] Lý thuyết về cấu tạo, tính chất vật lí của tinh bột và cellulose
-------------[2] Lý thuyết về tính chất hóa học (phản ứng thủy phân, phản ứng màu với iodine của tinh bột)
-------------[3] Viết các phương trình hóa học minh họa
-------------[4] Bài tập về ứng dụng của tinh bột và cellulose trong thực phẩm, công nghiệp
-------------[5] Bài toán thủy phân tinh bột/cellulose (tính hiệu suất)
-------------[6] Bài toán liên quan đến chuỗi phản ứng: Tinh bột → glucose → alcohol → acetic acid

----------[4] Protein
-------------[1] Lý thuyết về khái niệm, cấu tạo và vai trò của protein
-------------[2] Lý thuyết về tính chất của protein (đông tụ, thủy phân, phản ứng màu biuret)
-------------[3] Bài tập nhận biết protein và phân biệt với các chất hữu cơ khác
-------------[4] Bài tập liên quan đến vai trò của protein trong dinh dưỡng và đời sống

----------[5] Polymer
-------------[1] Lý thuyết về khái niệm, phân loại và tính chất chung của polymer
-------------[2] Lý thuyết về tính chất và ứng dụng của một số polymer thông dụng (PE, PVC, cao su, tơ)
-------------[3] Viết phương trình phản ứng trùng hợp điều chế polymer
-------------[4] Bài tập nhận biết và phân loại các loại vật liệu polymer (chất dẻo, tơ, cao su)
-------------[5] Bài toán tính khối lượng monomer/polymer (có hiệu suất phản ứng)
-------------[6] Bài tập về vấn đề ô nhiễm môi trường do rác thải nhựa và các giải pháp

-------[A] Khai thác tài nguyên từ vỏ Trái Đất

----------[1] Sơ lược về hoá học vỏ Trái Đất và khai thác tài nguyên từ vỏ Trái Đất

----------[2] Khai thác đá vôi. Công nghiệp silicate

----------[3] Khai thác nhiên liệu hoá thạch. Nguồn carbon. Chu trình carbon và sự ấm lên toàn cầu

-------[B] Di truyền học Mendel. Cơ sở phân tử của hiện tượng di truyền

----------[1] Khái quát về di truyền học

----------[2] Các quy luật di truyền của Mendel
-------------[1] Thí nghiệm lai một cặp tính trạng, quy luật phân li
-------------[2] Thí nghiệm lai hai cặp tính trạng, quy luật phân li độc lập
-------------[3] Bài tập xác định kiểu gen, kiểu hình, tỉ lệ phân li
-------------[4] Bài tập lai một tính trạng (xác định gen trội/lặn, viết sơ đồ lai)
-------------[5] Bài tập lai hai tính trạng
-------------[6] Liên hệ thực tiễn: giải thích tại sao con giống bố mẹ nhưng không hoàn toàn giống

----------[3] Nucleic acid và gene
-------------[1] Cấu trúc ADN (mô hình Watson-Crick)
-------------[2] Cơ chế tự nhân đôi ADN
-------------[3] Gen và mối quan hệ gen - ARN - Protein
-------------[4] Bài tập tính chiều dài, số nucleotide, liên kết hydro của ADN
-------------[5] Liên hệ: ứng dụng ADN trong xác định huyết thống, điều tra hình sự

----------[4] Tái bản DNA và phiên mã tạo RNA

----------[5] Dịch mã và mối quan hệ từ gene đến tính trạng

----------[6] Đột biến gene
-------------[1] Đột biến gen và đột biến NST (khái niệm, nguyên nhân)
-------------[2] Thường biến (khái niệm, đặc điểm)
-------------[3] Phân biệt biến dị di truyền và biến dị không di truyền
-------------[4] Liên hệ thực tiễn: tác nhân gây đột biến trong đời sống (thuốc trừ sâu, tia UV, hóa chất)

-------[C] Di truyền nhiễm sắc thể

----------[1] Nhiễm sắc thể và bộ nhiễm sắc thể
-------------[1] Cấu trúc nhiễm sắc thể, bộ NST của người (2n = 46)
-------------[2] Nguyên phân và giảm phân
-------------[3] Cơ chế xác định giới tính (NST X, Y)
-------------[4] Bài tập về NST trong nguyên phân, giảm phân
-------------[5] Giải thích tại sao tỉ lệ sinh con trai/gái xấp xỉ 1:1

----------[2] Nguyên phân và giảm phân

----------[3] Nhiễm sắc thể giới tính và cơ chế xác định giới tính

----------[4] Di truyền liên kết

----------[5] Đột biến nhiễm sắc thể

-------[D] Di truyền học với con người và đời sống

----------[1] Di truyền học với con người
-------------[1] Khái niệm quần thể, các đặc trưng (mật độ, tỉ lệ giới tính, thành phần tuổi)
-------------[2] Sự tăng trưởng quần thể
-------------[3] Liên hệ: vấn đề dân số Việt Nam và thế giới, kế hoạch hóa gia đình
-------------[4] Khái niệm quần xã, các mối quan hệ trong quần xã (cộng sinh, ký sinh, cạnh tranh, hội sinh)
-------------[5] Hệ sinh thái: thành phần, chuỗi thức ăn, lưới thức ăn
-------------[6] Phân tích hệ sinh thái tại địa phương (ao, vườn, rừng, đồng ruộng)
-------------[7] Giải thích hiện tượng mất cân bằng sinh thái khi một loài bị diệt (sâu bọ, rắn, ếch...)
-------------[8] Tác động của con người đến môi trường (ô nhiễm, phá rừng, biến đổi khí hậu)
-------------[9] Ô nhiễm môi trường: nguồn gây ô nhiễm, hậu quả, biện pháp khắc phục
-------------[A] Bảo vệ đa dạng sinh học và phát triển bền vững
-------------[B] Đánh giá thực trạng môi trường tại địa phương và đề xuất giải pháp
-------------[C] Xây dựng kế hoạch hành động bảo vệ môi trường tại trường học

----------[2] Ứng dụng công nghệ di truyền vào đời sống
-------------[1] Ứng dụng di truyền trong chọn giống (lai, gây đột biến, công nghệ gen)
-------------[2] Di truyền y học: bệnh di truyền ở người (Down, bạch tạng, máu khó đông)
-------------[3] Tư vấn di truyền và sàng lọc trước sinh
-------------[4] Thảo luận: GMO (thực phẩm biến đổi gen) - lợi ích và rủi ro

-------[E] Tiến hoá

----------[1] Khái niệm tiến hoá và các hình thức chọn lọc

----------[2] Cơ chế tiến hoá
-------------[1] Bằng chứng tiến hóa (hóa thạch, cơ quan tương đồng, thoái hóa)
-------------[2] Học thuyết tiến hóa của Darwin (chọn lọc tự nhiên)
-------------[3] Giải thích sự đa dạng sinh vật trên Trái Đất

----------[3] Sự phát sinh và phát triển sự sống trên Trái Đất
-------------[1] Các giai đoạn phát sinh sự sống trên Trái Đất
-------------[2] Sự phát sinh loài người
-------------[3] Liên hệ: bảo tồn đa dạng sinh học, bảo vệ các loài có nguy cơ tuyệt chủng
